
\documentclass[12]{article}
\usepackage{graphicx}
\usepackage[T1]{fontenc}
\usepackage{tgbonum}
\usepackage{amsmath}
\usepackage{amssymb}
\usepackage{physics}
\usepackage{booktabs}
\usepackage{float}
\usepackage{graphicx} % Required for inserting images
\usepackage{ulem}
%\usepackage{hyperref}
\usepackage{fontawesome5}
\usepackage{dirtytalk}
\usepackage{xcolor}
\usepackage[affil-it]{authblk}

\usepackage{setspace}
\doublespacing

\usepackage{geometry}
\geometry{
 a4paper,
 total={170mm,257mm},
 left=20mm,
 top=20mm,
 }

\usepackage[colorlinks=true,linkcolor=blue,urlcolor=black]{hyperref}
\usepackage{bookmark}




\title{EK369E: Financial Markets\\ Bond Prices and Yields\\ Problem Set (with solutions)}
\author{Nord University Business School}
\date{}

\begin{document}

\maketitle

\color{black}
\section*{Problem 1}
Consider a 6\% coupon bond selling for NOK 93.3 with three years to maturity making annual coupon payments. This bond is redeemable at maturity at par value.
\begin{enumerate}
    \item[a)] Calculate the bond’s yield to maturity (YTM).  
\end{enumerate}

\color{blue}
%\subsection*{Solution to Problem 1}
\subsubsection*{Solution 1a.}
Redemption/par value of the bond: NOK 100\\
Coupon payments on this bond: 6\% of 100 = NOK 6\\
Therefore, the YTM can be solved using the equation:\\
93.3 ={\Large $ \frac{6}{(1+YTM)^1} + \frac{6}{(1+YTM)^2} + \frac{100 +6}{(1+YTM)^3}$}\\


\noindent
\\
\textbf{Using the trial and error method}\\
%\subsubsection*{Using the trial and error method}
The YTM of this bond must be higher than 6\% since the bond is trading at a discount.
Given the difference between the bond's price and the par value, we can begin with a YTM of 7.5\% and increase incrementally until we get the price that equates the present value of inflows to the present value of out flows.\\
We can also speed up the process by using the annuity formula to discount the coupon payments of the bond. Annuity formula = $\frac{1 - (1+ YTM)^{-T}}{YTM}$\\
@YTM = 7.5\%\\
$P_0 = $ {\Large $= 6[\frac{1 - (1.075)^{-3}}{0.075}] + \frac{100}{(1.075)^3}$}\\
$P_0 = (6 \times 2.6005) + 80.4961 = 96.099$\\
This price is higher than 93.3, so we increase the YTM by 1\%\\
\\
YTM = 8.5\%\\
$P_0 = $ {\Large $= 6[\frac{1 - (1.085)^{-3}}{0.085}] + \frac{100}{(1.085)^3}$}\\
$P_0 = (6\times 2.5540) + 78.291 = 93.61$\\
This price is slightly higher than 93.3, so we increase the YTM by 0.05\%\\
\\
@YTM = 8.55\%\\
$P_0 = $ {\Large $= 6[\frac{1 - (1.0855)^{-3}}{0.0855}] + \frac{100}{(1.0855)^3}$}\\
$P_0 = (6\times 2.5517) + 78.1826 = 93.49$\\
This price is approximately the same as 93.3, we can conclude that the 
\uuline{YTM is approximately 8.55\%}

\noindent
\\
\textbf{Solving as a polynomial equation to the degree 3}\\
%\subsubsection*{Solving as a polynomial equation to the degree 3}
We set $(1 + YTM) = x$\\
{93.3 = \Large $\frac{6}{x^1} + \frac{6}{x^2} + \frac{106}{x^3}$}  \\
\\
{93.3 = \Large $\frac{6x^2 + 6x + 106}{x^3}$}  \hspace{4mm}  {\small(\textit{by finding the Least common multiple on the right side of the equation above})}\\
\\
{$93.3x^3 - 6x^2 -6x -106$} \hspace{4mm}  {\small(\textit{by cross multiplying (or multiply both sides of the the equation above by $x^3$) and equate to zero})}\\
Solving for x using equation mode on your calculator yields:\\
$ x = 1.086 \therefore \uuline{YTM = 0.086 \text{ or } 8.6\%}$

\noindent
\\
\textbf{Solving for YTM using a financial calculator}\\
%\subsubsection*{Solving Using a financial calculator}
\begin{table}[H]
    \centering
    \resizebox{0.5\textwidth}{!}{%
    {\color{blue}\begin{tabular}{lcl}
    \toprule
        Bond Information & Financial Calculator Input & Value\\\midrule
        Time to Maturity & N  & 3\\
        Par value & FV & 100\\
        Current price  & PV & -93.3\\
        Coupon payment & PMT & 6\\
        YTM & CPT + I/Y & \uuline{8.63}\\\bottomrule
    \end{tabular}}
    }
    \label{tab:my_label}
\end{table}
The financial calculator can solve for any other missing information about the bond by simply inputting the available bond information and pressing \say{CPT + missing information}. For example to solve for the current price of the bond, I input the information for N, FV, I/Y, PMT and press \say{CPT + PV}\\


\color{black}
\begin{enumerate}
    \item[b)] If interest rates over the next three years will be, with certainty, $r_1 = 6\%$, $r_2 = 8\%$ and $r_3 = 10\%$, calculate the realised compound average return on this bond.
\end{enumerate}


\color{blue}
\subsubsection*{Solution 1b}
\noindent
The compound rate or return takes into consideration income from reinvested coupon payments over three years plus capital gain and satisfies the relationship $V_0(1+r)^3 = V_3$\\
where $V_0$ is the invested amount, $V_3$ is the settlement amount at the end of year 3 and $r$ is the Geometric/Time Weighted/ Compound Return. 


\begin{table}[H]
    \centering
    \resizebox{0.8\textwidth}{!}{%
    {\color{blue}\begin{tabular}{cll}
    \toprule
         & Description & Cash flow \\\midrule
         & Ist coupon received at t = 1 and reinvested   & $6\times 1.08 \times1.1$ = 7.13\\
         & 2nd coupon received at t = 2 and reinvested & $6\times1.1$ = 6.6\\
         &3rd coupon received  & 6\\
         & Bond redemption  & 100 \\\bottomrule
         & Settlement amount  & 119.73 \\\bottomrule
    \end{tabular}}
    }
    \label{tab:my_label}
\end{table}

\noindent
Invested amount ($V_0$) is NOK 93.3\\
Settlement amount ($V_3$) is NOK 119.73\\
Geometric/Time Weighted/ Compound Return =  {\Large $\sqrt[3]{\frac{119.73}{93.3}} - 1$} = \uuline{8.7\%}\\
Note that the compound rate of return is not equal to the YTM ($8.7 \neq 8.6$) since coupon was not re-invested at the YTM.\\
\\
If the coupon received was reinvested at the YTM of 8.6\%
\begin{table}[H]
    \centering
    \resizebox{0.8\textwidth}{!}{%
    {\color{blue}\begin{tabular}{cll}
    \toprule
         & Description & Cashflow \\\midrule
        & Ist coupon received at t = 1and reinvested   & $6\times 1.086 \times1.086$ = 7.07\\
         & 2nd coupon received at t = 2 and reinvested & $6\times1.086$ = 6.52\\
         &3rd coupon received at t = 3  & 6\\
         & Bond redemption  & 100 \\\bottomrule
         & Settlement amount  &  119.59\\\bottomrule
    \end{tabular}}
    }
    \label{tab:my_label}
\end{table}

\noindent
Invested amount ($V_0$) is NOK 93.3\\
Settlement amount ($V_3$) is NOK 119.59\\
Geometric/Time Weighted/ Compound Return =  {\Large $\sqrt[3]{\frac{119.59}{93.3}} - 1$} = \uuline{8.6\%}\\







\color{black}
\section*{Problem 2}
A 20 Year maturity bond with par value of NOK 100 makes annual coupon payments at 8\%. Find the yield to maturity of the bond if the bond's price is
\begin{enumerate}
    \item[a)] NOK 95
    \item[b)] NOK 100
    \item[c)] NOK 105
\end{enumerate}

\color{blue}
\subsection*{Solution to Problem 2}
\subsubsection*{Solution 2a.}
Before making computations, we know that a bond trading at a discount must have its YTM greater than its coupon rate. We can verify this by computing YTM using the financial calculator.
\begin{table}[H]
    \centering
    \resizebox{0.5\textwidth}{!}{%
    {\color{blue}\begin{tabular}{lcl}
    \toprule
        Bond Information & Financial Calculator Input & Value\\\midrule
        Time to Maturity & N  & 20\\
        Par value & FV & 100\\
        Current price  & PV & -95\\
        Coupon payment & PMT & 8\\
        YTM & CPT + I/Y & \uuline{8.53\%}\\\bottomrule
    \end{tabular}}
    }
    \label{tab:my_label}
\end{table}


\subsubsection*{Solution 2b.}
Without computations, we know that a bond trading at par has its YTM equal to its coupon rate. We can verify this by computing YTM using the financial calculator.
\begin{table}[H]
    \centering
    \resizebox{0.5\textwidth}{!}{%
    {\color{blue}\begin{tabular}{lcl}
    \toprule
        Bond Information & Financial Calculator Input & Value\\\midrule
        Time to Maturity & N  & 20\\
        Par value & FV & 100\\
        Current price  & PV & -100\\
        Coupon payment & PMT & 8\\
        YTM & CPT + I/Y & \uuline{8.00\%}\\\bottomrule
    \end{tabular}}
    }
    \label{tab:my_label}
\end{table}


\subsubsection*{Solution 2c.}
Before making computations, we know that a bond trading at a premium must have its YTM less than its coupon rate. We can verify this by computing YTM using the financial calculator.
\begin{table}[H]
    \centering
    \resizebox{0.5\textwidth}{!}{%
    {\color{blue}\begin{tabular}{lcl}
    \toprule
        Bond Information & Financial Calculator Input & Value\\\midrule
        Time to Maturity & N  & 20\\
        Par value & FV & 100\\
        Current price  & PV & -105\\
        Coupon payment & PMT & 8\\
        YTM & CPT + I/Y & \uuline{7.51\%}\\\bottomrule
    \end{tabular}}
    }
    \label{tab:my_label}
\end{table}


\color{black}
\section*{Problem 3}
Fill in the table below for the following zero-coupon bonds with par values of NOK 100.

\begin{table}[H]
    \centering
    \resizebox{0.5\textwidth}{!}{%
    \begin{tabular}{lcc}
    \toprule
        Price & Years to Maturity & Yield To Maturity (YTM)\\\midrule
        NOK 40 & 20 & a) \\
        NOK 50 & 20 & b)\\
        NOK 50 & 10 & c)\\
        d) & 10 & 10\%\\
        e) & 10 & 8\%\\
        NOK 40 & f) & 8\%\\\bottomrule
    \end{tabular}
    }
    \label{tab:my_label}
\end{table}


\color{blue}
\subsection*{Solution to Problem 3}
\subsubsection*{Solution 3a.}
\begin{table}[H]
    \centering
    \resizebox{0.5\textwidth}{!}{%
    {\color{blue}\begin{tabular}{lcl}
    \toprule
        Bond Information & Financial Calculator Input & Value\\\midrule
        Time to Maturity & N  & 20\\
        Par value & FV & 100\\
        Current price  & PV & -40\\
        Coupon payment & PMT & 0\\
        YTM & CPT + I/Y & \uuline{4.69\%}\\\bottomrule
    \end{tabular}}
    }
    \label{tab:my_label}
\end{table}

\subsubsection*{Solution 3b.}
\begin{table}[H]
    \centering
    \resizebox{0.5\textwidth}{!}{%
    {\color{blue}\begin{tabular}{lcl}
    \toprule
        Bond Information & Financial Calculator Input & Value\\\midrule
        Time to Maturity & N  & 20\\
        Par value & FV & 100\\
        Current price  & PV & -50\\
        Coupon payment & PMT & 0\\
        YTM & CPT + I/Y & \uuline{3.53\%}\\\bottomrule
    \end{tabular}}
    }
    \label{tab:my_label}
\end{table}


\subsubsection*{Solution 3c.}
\begin{table}[H]
    \centering
    \resizebox{0.5\textwidth}{!}{%
    {\color{blue}\begin{tabular}{lcl}
    \toprule
        Bond Information & Financial Calculator Input & Value\\\midrule
        Time to Maturity & N  & 10\\
        Par value & FV & 100\\
        Current price  & PV & -50\\
        Coupon payment & PMT & 0\\
        YTM & CPT + I/Y & \uuline{7.18\%}\\\bottomrule
    \end{tabular}}
    }
    \label{tab:my_label}
\end{table}



\subsubsection*{Solution 3d.}
The price of the bond is $P_0 = ${\Large $ \frac{100}{(1.1)^{10}}$} =   \uuline{NOK 38.55}
\begin{table}[H]
    \centering
    \resizebox{0.5\textwidth}{!}{%
    {\color{blue}\begin{tabular}{lcl}
    \toprule
        Bond Information & Financial Calculator Input & Value\\\midrule
        Time to Maturity & N  & 10\\
        Par value & FV & 100\\
        Coupon payment & PMT & 0\\
        YTM & I/Y & 10\%\\
        Current price  & CPT + PV & \uuline{-38.55}\\\bottomrule
    \end{tabular}}
    }
    \label{tab:my_label}
\end{table}
\noindent
%Or \\


\subsubsection*{Solution 3e.}
The price of the bond is $P_0 = ${\Large $ \frac{100}{(1.08)^{10}}$} = \uuline{NOK 46.32}
\begin{table}[H]
    \centering
    \resizebox{0.5\textwidth}{!}{%
    {\color{blue}\begin{tabular}{lcl}
    \toprule
        Bond Information & Financial Calculator Input & Value\\\midrule
        Time to Maturity & N  & 10\\
        Par value & FV & 100\\
        Coupon payment & PMT & 0\\
        YTM & I/Y & 8\%\\
        Current price  & CPT + PV & \uuline{-46.32}\\\bottomrule
    \end{tabular}}
    }
    \label{tab:my_label}
\end{table}
%\noindent
%Or\\
%\noindent



\subsubsection*{Solution 3f.}
$40 = $ {\Large $\frac{100}{1.08^T}$}\\
$1.08^T =${\Large $\frac{100}{40}$}\\
$\log(1.08^T) = \log(2.5)$\\
$T\cdot \log(1.08) = \log(2.5)$\\
$T = \frac{\log(2.5)}{\log(1.08)} = \frac{0.3979}{0.0334} \approx 11.9 $\\
\noindent
The time to maturity of the bond is approximately \uuline{11.9 years}

\noindent
\begin{table}[H]
    \centering
    \resizebox{0.5\textwidth}{!}{%
    {\color{blue}\begin{tabular}{lcl}
    \toprule
        Bond Information & Financial Calculator Input & Value\\\midrule
        Par value & FV & 100\\
        Coupon payment & PMT & 0\\
        YTM & I/Y & 8\%\\
        Current price  & PV & - 40\\
        Time to Maturity & CPT + N  & \uuline{11.9}\\\bottomrule
    \end{tabular}}
    }
    \label{tab:my_label}
\end{table}

\noindent





\color{black}
\section*{Problem 4}
Bond Alfa has a current yield of 9\% and a yield to maturity of 10\%. 
\begin{enumerate}
    \item[a)] Is Bond Alfa selling above or below par value?   
\end{enumerate}

\color{blue}
\subsubsection*{Solution 4a}
A bond selling above par value i.e. a premium bond has a yield to maturity less than its current yield, which is also less than its coupon rate\\
A bond selling below par value i.e. a discount bond has a yield to maturity greater than its current yield, which is also greater than its coupon rate\\
Bond Alfa has a yield to maturity greater than its current yield therefore, Bond Alfa is trading at a discount (below par value).\\


\color{black}
\begin{enumerate}
    \item[b)] Is the coupon rate on Bond Alfa more or less than 9\%?
\end{enumerate}

\color{blue}
\subsubsection*{Solution 4b}
Since Bond Alfa is trading at a discount, its coupon rate is less than 9\% (i.e. its coupon rate is less than its current yield and its yield to maturity).



\color{black}
\section*{Problem 5}
A newly issued bond pays its coupons once annually. Its coupon rate is 5\%, its maturity is 20 years and its yield to maturity is 8\%.
\begin{enumerate}
    \item[a)] Find the holding period return for a 1-year investment period if the bond is selling at a yield to maturity of 7\% by the end of the first year.
\end{enumerate}
\color{blue}
\subsubsection*{Solution 5a.}
Price at the beginning of the year ($P_0$):\\
$P_0 = $ {\Large $= 5[\frac{1 - (1.08)^{-20}}{0.08}] + \frac{100}{(1.08)^{20}}$} = 70.55\\
\\
$P_1 = $ {\Large $= 5[\frac{1 - (1.07)^{-19}}{0.07}] + \frac{100}{(1.07)^{19}}$} = 79.33\\
\\
$HPR_1$ = {\Large $\frac{79.33 - 70.55 + 5}{70.55}$ }= \uuline{0.1953 / 19.53\%}\\


\color{black}
\begin{enumerate}
    \item[b)] Continuing from a), find the realised compound yield for a two-year holding period, assuming that you sell the bond after two years, the bond's yield is 7\% at the end of the second year and coupons can be reinvested for one year at 3\% interest rate.
\end{enumerate}

\color{blue}
\subsubsection*{Solution 5b.}

$HPR_1 = 19.53\%$\\
\\
$P_2 = $ {\Large $= 5[\frac{1 - (1.07)^{-18}}{0.07}] + \frac{100}{(1.07)^{18}}$} = 79.88\\
\\
$HPR_2$ = {\Large $\frac{79.88 - 79.33 + 5}{79.33}$ }= \uuline{6.996\% or 7\%}\\
\\
Realised Compound Yield = $(\sqrt[2]{1.1953 \times 1.07}) - 1$ = \uuline{13.09\% $\approx$ 13\%}

\noindent
or 

\noindent


\begin{table}[H]
    \centering
    \resizebox{0.8\textwidth}{!}{%
    {\color{blue}\begin{tabular}{cll}
    \toprule
         & Description  & Cashflow\\\midrule
         & 1st coupon received at t = 1 and reinvested at 3\% & $5\times1.03$ = 5.15\\
         & 2nd  coupon received at t = 2  & 5\\
         & Proceeds from the sale of the bond & 79.88  \\\bottomrule
         & Settlement amount  & 90.03 \\\bottomrule
    \end{tabular}}
    }
    \label{tab:my_label}
\end{table}
\noindent
Invested amount ($V_0$) is NOK 70.55\\
Settlement amount ($V_2$) is NOK 90.03\\
Realised Compound Yield =  {\Large $\sqrt[2]{\frac{90.03}{70.55}} - 1$} =  \uuline{12.965\% $\approx$ 13\%}\\




\color{black}
\section*{Problem 6}
The table below summarises some information from the fixed-income market:
% Please add the following required packages to your document preamble:
% \usepackage{booktabs}
% \usepackage{graphicx}
\begin{table}[H]
\centering
\resizebox{0.4\textwidth}{!}{%
\begin{tabular}{@{}llll@{}}
\toprule
Bond & \begin{tabular}[c]{@{}l@{}}Maturity\\ (years)\end{tabular} & \begin{tabular}[c]{@{}l@{}}Coupon rate\\ (\%)\end{tabular} & \begin{tabular}[c]{@{}l@{}}Yield\\ (\%)\end{tabular} \\ \midrule
A    & 5.5                                                        & 0                                                          & 2.5                                                  \\
B    & 2.5                                                        & 4.0                                                        & 2.0                                                  \\ \bottomrule
\end{tabular}%
}
\label{tab:my-table}
\end{table}

\begin{enumerate}
    \item[a)] Calculate the invoice price and the flat price for the two bonds.
\end{enumerate}

\color{blue}
\subsubsection*{Solution 6a.}
Invoice price = Clean Price + Accrued interest (since last coupon payment)\\
\textbf{Bond A}\\
Bond A does not pay coupon , therefore its flat price is the same as its invoice price.\\
$P_{0, Clean}^A = P_{0, Invoice }^A = $ {\Large $\frac{100}{(1.025)^{5.5}}$} = \uuline{87.30}\\

\noindent
\textbf{Bond B}\\
There are two methods of computing the  invoice price.\\
\textit{Method 1}: compute the clean price based on T = 2.5 and add accrued coupon for 0.5 years to get invoice price\\
$P_{0, clean}^B = $ {\Large $= 4[\frac{1 - (1.020)^{-2.5}}{0.020}] + \frac{100}{(1.020)^{2.5}}$ }\\
$P_{0, clean}^B = $ {\large $= (4\times 2.42) + 95.17$} = \uuline{104.85}\\
Accrued interest on Bond B: $0.5\times 4$ = 2\\
$P_{0, invoice}^B$ =  104.85 + 2 = \uuline{106.85}\\

\noindent
\textit{Method 2}:
Compute the invoice price using the formula and setting T = 3\\
$P_{0, invoice}^B $ = {\huge $\frac{\frac{4}{0.02} \times [1.02^3 -1] + 100 }{(1.02)^{3-1+0.5}}$}\\
\\
$P_{0, invoice}^B $ = {\Large $\frac{112.24}{1.0508}$}  = \uuline{106.81}\\
\\
Accrued interest on Bond B: $0.5\times 4$ = 2\\
$P_{0, clean}^B =$ 106.81 - 2 = \uuline{104.81}



\noindent
\color{black}
\begin{enumerate}
    \item[b)] How can a zero-coupon bond be traded above its face value?
\end{enumerate}

\color{blue}
\subsubsection*{Solution 6b.}
A bond trades above its par value when its yield to maturity is less than its coupon rate. For a zero-coupon bond to trade at a premium, its yield to maturity must be less than zero i.e market interest rates must be negative.








\color{black}
\section*{Problem 7}
Treasury bonds paying an 8\% coupon rate with semiannual payments currently sell at par value. what coupon rate would they have to pay in order to sell at par if they paid their coupons annually.

\color{blue}
\subsubsection*{Solution 7}
A semi-annual coupon bond selling at par has semi-annual coupon rate = semi-annual YTM = 8\%\\
This represents an effective annual yield of $0.08^2 = 0.1664$ or $16.64\%$\\
Therefore,to continue to sell at par, when paying coupons annually, the bond's coupon rate must be equal to its YTM of \uuline{16.64\%}.





\color{black}
\section*{Problem 8}
Assume that you have a 1-year investment horizon and are trying to choose among three bonds with the same degree of default risk and the same time to maturity of 10 years. Bond 1 is a zero coupon bond, Bond 2 has a coupon rate of 8\%, and Bond 3 has a coupon rate of 10\%
\begin{enumerate}
    \item[a)] If all bonds are currently priced to yield 8\% at maturity, what are the prices of Bond 1, Bond 2 and Bond 3
\end{enumerate}


\color{blue}
\subsubsection*{Solution 8a.}
$P_0^{Bond 1} =$ {\Large $\frac{100}{1.08^{10}}$} = \uuline{ 46.32 }\\
\\
$P_0^{Bond 2} =$ \uuline{100}\\
\\
$P_0^{Bond 3} =$ {\Large $= 10[\frac{1 - (1.08)^{-10}}{0.08}] + \frac{100}{(1.08)^{10}}$} = \uuline{113.42} \\


\noindent
\color{black}
\begin{enumerate}
    \item[b)] If at the beginning of next year, all bonds are priced to yield 8\% , what will the prices of these bonds be?
\end{enumerate}

\color{blue}
\subsubsection*{Solution 8b.}
$P_1^{Bond 1} =$ {\Large $\frac{100}{1.08^{9}}$} = \uuline{50.02}\\
\\
$P_1^{Bond 2} =$ \uuline{100}\\
\\
$P_1^{Bond 3} =$ {\Large $= 10[\frac{1 - (1.08)^{-9}}{0.08}] + \frac{100}{(1.08)^{9}}$} = \uuline{112.49} \\

\noindent
Comments: \\
Holding all else equal and assuming stable interest rates, discount bonds increase in value (approach the par-value) as they approach maturity. \\
Holding all else equal and assuming stable interest rates, premium bonds reduce in value (approach the par-value) as they approach maturity. This is because the number of high coupon payments, the reason why they were priced above par, reduces as maturity approaches.\\
In summary, bonds(with similar attributes) are priced such that all investors earn comparable total return (otherwise there would be arbitrage in the bond market). 



\noindent
\color{black}
\begin{enumerate}
    \item[c)] Based on your answer in a) and b) compute the holding period return over the 1 year horizon for each bond. 
\end{enumerate}


\color{blue}
\subsubsection*{Solution 8c.}
$HPR_1^{Bond 1}$ = {\Large $\frac{50.02 - 46.32 }{46.32}$ }= \uuline{7.99\% $\approx$ 8.0\%}\\
\\
$HPR_1^{Bond 2}$ = {\Large $\frac{100 - 100 + 8 }{100}$ }= \uuline{8.0\%}\\
\\
$HPR_1^{Bond 3}$ = {\Large $\frac{112.49 - 113.42 + 10 }{113.42}$ }= \uuline{7.997\% $\approx$ 8.0\%}\\

\noindent
Comments: 
With stable interest rates, these bonds are interchangeable as the investor earns the same return regardless.




\noindent
\color{black}
\begin{enumerate}
    \item[d)] Recalculate b) and c) if at the beginning of next year, all bonds are priced to yield 7\%
\end{enumerate}

\color{blue}
\subsubsection*{Solution 8d.}
$P_1^{Bond 1} =$ {\Large $\frac{100}{1.07^{9}}$} = \uuline{54.39}\\
\\
$P_1^{Bond 2} =${\Large $= 8[\frac{1 - (1.07)^{-9}}{0.08}] + \frac{100}{(1.07)^{9}}$} =  \uuline{106.52}\\
\\
$P_1^{Bond 3} =$ {\Large $= 10[\frac{1 - (1.07)^{-9}}{0.08}] + \frac{100}{(1.07)^{9}}$} = \uuline{119.55} \\
\\
\\
$HPR_1^{Bond 1}$ = {\Large $\frac{54.39 - 46.32 }{46.32}$ }= \uuline{17.42\% }\\
\\
$HPR_1^{Bond 2}$ = {\Large $\frac{106.52 - 100 + 8 }{100}$ }= \uuline{14.52\%}\\
\\
$HPR_1^{Bond 3}$ = {\Large $\frac{119.55 - 113.42 + 10 }{113.42}$ }= \uuline{14.22\%}\\

\noindent
Comment: 
When interest rates decline, the HPR of a bond increases (due to the capital appreciation in price). This appreciation is more for zero-coupon bonds (they more quickly approach par-value). In this example, the zero-coupon bond performs better than the coupon-paying bonds since majority of the holding period return comes from capital appreciation when interest rates reduce.
\end{document}


